\chapter{Cryptographic Provenance and the Security Model}
\label{ch:security}

This chapter makes explicit the cryptography that underlies every claim of
immutability in this monograph. It also seals the manuscript itself.

\section{The Three Anchors}
\label{sec:sec-anchors}

Three cryptographic primitives anchor the system (and this paper):
\begin{enumerate}
  \item \textbf{SHA-256 digital fingerprint} --- a content hash that
        identifies an artifact uniquely.
  \item \textbf{Ed25519 signature} --- an author-attributed, verifiable
        attestation over the fingerprint.
  \item \textbf{Bitcoin OP\_RETURN} --- an economic, adversarial-resistant
        timestamp that anchors the fingerprint into the Bitcoin ledger.
\end{enumerate}

\section{Digital Fingerprint (SHA-256)}
\label{sec:sec-sha}

The fingerprint of this manuscript's source tree is computed over the
concatenation of its chapter files and master document:
\begin{center}
\ttfamily\small \input{anchors/fingerprint.txt}
\end{center}
Any reader may recompute it:
\begin{verbatim}
cat mathlib5_full.tex chapters/*.tex | sha256sum
\end{verbatim}
If the value differs, the manuscript has been altered since sealing.

\section{Ed25519 Attestation}
\label{sec:sec-ed}

The fingerprint is signed with an Ed25519 private key whose public component
is:
\begin{center}
\ttfamily\small \input{anchors/pubkey.txt}
\end{center}
The signature is recorded in \texttt{anchors/signature.txt} and can be
verified with any Ed25519 implementation:
\begin{verbatim}
python -c "import cryptography; ..."  # verify fingerprint against pubkey
\end{verbatim}
The signature answers the question \emph{who authored this?} The Plasma Gate
of \cref{sec:coh-plasma} enforces exactly this check on every artifact.

\section{Bitcoin OP\_RETURN Anchor}
\label{sec:sec-btc}

To make the fingerprint durable against even the compromise of the author's
infrastructure, it is anchored into Bitcoin via an \texttt{OP\_RETURN}
output. The anchor message is:
\begin{center}
\ttfamily\footnotesize \input{anchors/opreturn.txt}
\end{center}
An \texttt{OP\_RETURN} is a provably-unspendable output carrying up to 80
bytes of arbitrary data. Once a transaction containing it is mined, the
fingerprint is timestamped by the Bitcoin consensus at a known block height,
with economic finality: reversing it would require re-organizing the chain,
costing millions of dollars. The anchor thus provides \emph{external,
trust-minimized} provenance that no single party controls.

\section{How to Verify This Manuscript}
\label{sec:sec-verify}

\begin{enumerate}
  \item Compute the SHA-256 of the source tree; compare to
        \texttt{anchors/fingerprint.txt}.
  \item Verify the Ed25519 signature over that fingerprint using
        \texttt{anchors/pubkey.txt}.
  \item Look up the Bitcoin transaction whose \texttt{OP\_RETURN} contains the
        anchor message; confirm the block height and timestamp.
\end{enumerate}
All three must hold for the manuscript to be considered sealed. The scripts
in \texttt{anchors/} automate steps 1--2.

\section{The WORM + Signature + Encryption Stack}
\label{sec:sec-stack}

The full security posture of a \textsc{mathlib5} artifact is therefore a stack:
\begin{enumerate}
  \item \textbf{Integrity}: SHA-256 chain (append-only) --- tamper evidence.
  \item \textbf{Confidentiality}: AES-256-GCM --- encryption at rest.
  \item \textbf{Authenticity}: Ed25519 --- the Plasma Gate.
  \item \textbf{Durability}: Bitcoin OP\_RETURN --- external timestamp.
\end{enumerate}
This is the realization of the title's promise: an \emph{immutable system of
safety boundaries}. Each layer of the stack is independently verifiable and,
together, they make undetected modification of a sealed artifact
computationally and economically infeasible.

\section{Key Management}
\label{sec:sec-keys}

The Ed25519 Plasma-Gate keypair is the root of all boundary signatures.
Its public key is published (see \texttt{pubkey.txt}); the private key
must live only in the signing environment and never in the repository. The
build gate refuses to sign with a key that is not isolated, so a leaked
key cannot silently re-sign artifacts. Rotation is supported by anchoring
a new public key in the Bifrost genesis successor block; old signatures
remain valid for the interval they covered.
